% ==========================================================
% titles/title-02.tex -- Title II: Governance Design and Sortition
% ==========================================================

\ConstTitle{Governance Design and Sortition}{title:sortition}

\Article{Principles and Template of Governance Design}{art:governance-principles}

\ArtSection{Principles}{art:governance-principles:s-principles}
The governance of the Commonwealth is designed around the following principles:
\begin{itemize}
  \item No single person, body, or institution holds authority across the full scope of governance
  \item All governance roles are temporary, non-renewable, and post-service constrained, except for technically specialized positions of narrow scope that are subject to the reconfirmation mechanism in \ref{art:governance-principles:s-renewability}
  \item Random selection eliminates the self-selection of the power-hungry that elections produce
  \item Transparency makes informal power accumulation visible before it becomes entrenched
  \item Decisions are made at the level closest to the people they affect
  \item Redundancy and incoherence prevent single-point capture
  \item No governance body may set the rules that govern its own composition, selection, accountability, or the scope of its own authority. These functions are always assigned to an external body with independent democratic legitimacy. This principle is a corollary of the anti-entrenchment principle and applies to all layers of governance without exception.
\end{itemize}

\ArtSection{Renewability Guide}{art:governance-principles:s-renewability}
Non-renewability is constitutionally mandated for all positions with broad policy-making authority, cross-domain jurisdiction, or significant public discretion. Positions that are narrow in scope, technically specialized, fully transparent, and subordinate to sortition oversight may permit renewal, subject to conflict of interest and sector exclusion provisions applicable to all governance roles. The specific renewal terms for each position are specified in the relevant article.

For positions designated as renewable, a reconfirmation mechanism applies. The relevant sortition oversight body conducts a public review of the position holder's performance against constitutional mandates at the end of each term, or every five years if the term is longer than five years, unless a different interval is specified in the relevant article. Findings are published through the \gls{gls:transparency-engine}. The relevant \gls{gls:citizen-oversight-panels} submits an independent assessment. Reconfirmation requires a supermajority of the oversight body. Non-confirmation triggers a new \gls{gls:qualification-commons} selection process; the outgoing holder remains in post for up to ninety days during transition. Reconfirmation review is triggered automatically on a calendar basis. A \gls{gls:citizen-oversight-panels} supermajority may trigger early reconfirmation review at any time. Three percent of civic membership may trigger early reconfirmation review through the \gls{gls:civic-initiative-infrastructure} under Article \ref{art:civic-initiative-infrastructure}, subject to an eighteen-month cooling period per position.

\ArtSection{Standard Commons Bodies}{art:governance-principles:s-standard-commons-bodies}
Unless an article specifies otherwise, any body designated a standard body is:
\begin{itemize}
  \item selected by sortition from the relevant \gls{gls:qualification-commons} pool or the general civic membership
  \item subject to staggered, non-renewable terms
  \item accompanied by an embedded \gls{gls:citizen-oversight-panels} under Article \ref{art:citizen-oversight-panels}
  \item has its proceedings fully logged through the \gls{gls:transparency-engine} under Article \ref{art:transparency-engine}
\end{itemize}

\ArtSection{Conflict of Interest and Sector Exclusion}{art:governance-principles:s-conflict-interest-sector}
Before assuming any governance role, a person undergoes full financial and relational disclosure. During service in any governance role, a person may not receive income beyond their public salary and any applicable civic service premium, communicate privately with any entity that has interests in a matter pending before their governance body, or hold productive assets beyond personal property.

No person may deliberate on a matter in which they stand to gain or lose materially from a specific outcome in ways distinguishable from the general population. This exclusion applies to ownership and financial interest relationships, not to labor relationships. The \gls{gls:citizen-oversight-panels} embedded within the relevant governance body determines conflict of interest questions on a case-by-case basis, with full transparency, subject to \gls{gls:constitutional-court} review.

\Article{Designing Sortition}{art:sortition-design}

\ArtSection{What Sortition Is}{art:sortition-design:s-what-sortition}
Sortition is the selection of governance participants by random lottery, weighted to produce demographic representation across the civic membership. It replaces election as the primary mechanism for filling governance roles because elections systematically select for the desire for power -- precisely the trait that governance design must exclude.

\ArtSection{Mandatory Service, Civic Duty, and Hardship Exemptions}{art:sortition-design:s-mandatory-service-civic}
Selection to any national governance role is a mandatory civic duty, considered non-arbitrary coercion to facilitate a functioning state. The Constitution aims to make this selection desirable through civic premiums, but acknowledges that making the switch from worker to legislator may come with stress. Therefore, hardship exemptions are available on the following constitutional grounds:
\begin{itemize}
  \item \textbf{Medical incapacity} -- a condition preventing meaningful participation in governance duties, certified by a Commonwealth healthcare provider with no financial interest in the outcome.
  \item \textbf{Sole caregiver status} -- the selected person is the sole provider of care for a dependent whose welfare would be seriously endangered by their absence, with no available alternative caregiver.
  \item \textbf{Genuinely irreplaceable critical role} -- the selected person holds a role whose immediate vacancy would create serious risk to public safety, assessed by the relevant \gls{gls:implementation-councils} against published criteria. The Critical Role Phased Entry provision of \ref{art:sortition-design:s-critical-role-phased} must be exhausted before this exemption is granted.
  \item \textbf{Acute personal crisis} -- the selected person is experiencing a personal crisis of a severity that would prevent meaningful participation. This exemption is duration-limited; the person returns to the eligible pool once the crisis has passed.
\end{itemize}

Exemptions are self-certified against these categories and submitted to the \gls{gls:electoral-commons} for administrative verification. The \gls{gls:electoral-commons} makes no discretionary determination in routine cases -- verification confirms only that the certification matches a constitutional category. Disputed certifications are reviewed by the \gls{gls:exemption-review-panel}, a panel of nine persons selected by sortition within the \gls{gls:electoral-commons}, serving staggered two-year renewable terms subject to the reconfirmation mechanism of Article \ref{art:governance-principles}, \ref{art:governance-principles:s-renewability}. Decisions are logged in full through the \gls{gls:transparency-engine} and are appealable to the \gls{gls:constitutional-court} on constitutional grounds only. The Panel may not expand or narrow the constitutional exemption categories; it determines only whether a specific certification genuinely meets an existing category. The \gls{gls:electoral-commons} may develop subsidiary exemption criteria through public deliberation with mandatory community input, subject to \gls{gls:constitutional-court} review and \gls{gls:sortition-legislature} ratification. Exempted persons are returned to the eligible pool for future cycles unless the exemption ground is permanent.

\ArtSection{Service Protections and the Civic Service Premium}{art:sortition-design:s-service-protections-civic}
Serving members receive:
\begin{itemize}
  \item Their full pre-service salary, replaced by the Commons at no cost to the member
  \item A civic service premium of twenty-five percent above pre-service salary
  \item Guaranteed return to their previous position upon completion of service
  \item Continued accrual of seniority, pension rights, and professional standing during service
  \item An extended paid reintegration sabbatical upon return to civilian life
  \item Permanent civic service recognition
\end{itemize}

\ArtSection{Critical Role Phased Entry}{art:sortition-design:s-critical-role-phased}
For roles certified by relevant \gls{gls:implementation-councils}s as genuinely critical to immediate public safety, selected persons may negotiate a phased entry period of up to ninety days to transition their critical caseload before fully beginning service. This accommodation applies to the transition period only; it does not exempt persons from service.

\ArtSection{The Electoral Commons}{art:sortition-design:s-electoral-commons}
The \gls{gls:electoral-commons} is the Commons institution responsible for maintaining the civic membership roll, conducting all sortition, verifying eligibility, managing stratification methodology, and publishing full documentation of every selection process. It is governed by a standard body of twelve members serving two-year terms. Its selection algorithm is open-source and publicly auditable. The \gls{gls:electoral-commons} may not set rules governing its own composition, selection, or accountability; those functions are governed by this Constitution and the \gls{gls:sortition-legislature}.

\ArtSection{Stratification}{art:sortition-design:s-stratification}
The civic membership is divided into demographic strata -- by region, age range, gender, primary occupation category, and other characteristics determined by the \gls{gls:electoral-commons} through a public deliberative process with mandatory community input, subject to \gls{gls:constitutional-court} review and \gls{gls:sortition-legislature} ratification. The \gls{gls:electoral-commons} may not unilaterally expand or alter stratification dimensions. All stratification methodology is published in full through the \gls{gls:transparency-engine}.

\ArtSection{The Eligible Pool}{art:sortition-design:s-eligible-pool}
Any person who has achieved full civic membership under Article \ref{art:civic-integration} is eligible for selection by sortition unless they have previously served in the specific role being filled, hold a current conflict of interest in the relevant sector, or have been removed from a prior governance role for constitutional violation.

\ArtSection{The Qualification Commons}{art:sortition-design:s-qualification-commons}
Qualification pools for all sortition bodies requiring demonstrated competence are maintained by the \gls{gls:qualification-commons}, a permanent Commons institution. Qualification is determined by blind evaluation of demonstrated work, assessed against publicly published rubrics. Credentials, titles, and institutional affiliations are neither required nor considered. Rubrics are designed and revised by the \gls{gls:rubric-body}, a standard body of nine members serving three-year terms, stratified to ensure representation across all major qualification domains maintained by the \gls{gls:qualification-commons}. \gls{gls:rubric-body} members recuse from rubric design in any domain in which they hold professional standing; the remaining members constitute a valid quorum for that decision. \gls{gls:rubric-body} decisions require supermajority approval and a mandatory public comment period. Scoring panels are drawn by sortition from current and former members of the relevant pool and from qualified members of adjacent domains. All rubrics, scores, and processes are logged in full by the \gls{gls:transparency-engine}. Persons denied qualification may appeal to the \gls{gls:constitutional-court} solely on the ground that the rubric itself violated foundational principles. See Article \ref{art:founding-mechanisms}, \ref{art:founding-mechanisms:s-qualification-commons-bootstrap} for bootstrapping arrangements.

\ArtSection{Bootstrapping the First Legislature}{art:sortition-design:s-bootstrapping-first-legislature}
At founding, all sortition bodies are populated simultaneously. To establish staggered rotation, founding selections assign varying initial term lengths by lottery. For instance, one third of founding \gls{gls:sortition-legislature} members are assigned one-year founding terms, one third two-year terms, one third full three-year terms. Service in any founding term -- regardless of length -- triggers permanent non-renewability. Members assigned shortened founding terms are not financially disadvantaged; the civic service premium and post-service benefits are calculated on a pro-rated basis proportional to the term length served.

\ArtSection{Universal Onboarding}{art:sortition-design:s-universal-onboarding}
Every person selected by sortition to any governance body or \gls{gls:implementation-councils} receives, prior to assuming any voting or decision-making role, a mandatory orientation covering the foundational principles of this Constitution and the specific mandate and constraints of the body to which they are selected. Orientation content is designed by the \gls{gls:qualification-commons}, published in full through the \gls{gls:transparency-engine}, and required regardless of the pathway by which the selected person achieved civic membership. Orientation is educational only and may not be used to test, screen, or disqualify any selected person.